%% kapa: exam=1 round=23 year=2012 period=1 form=A subject=economics subject_ko=경제원론 seq=32 number=72
%% source: Q-Net 감정평가사 기출문제 (cst00309) artlSeq=5052578 file=raw/1008/5052578/2012년 제23회 감정평가사 제1차시험 기출문제(1교시 A형).pdf
%% conversion: kordoc 4.24.51 -> tools/gen1.py (auto); review=proofread-sonnet; flags=none
\begin{question}{72}
소비이론에 관한 설명으로 옳은 것은?
\begin{choices}
  \item 피셔(I. Fisher)의 기간 간 소비선택이론에 따르면 차입제약이 없는 경우 이자율은 현재소비에 영향을 줄 수 없다.
  \item 항상소득가설(permanent income hypothesis)은 소비자들이 유동성제약에 처해있다고 전제한다.
  \item 생애주기가설(life cycle hypothesis)은 현재소비는 현재소득에만 의존한다고 전제한다.
  \item 항상소득가설에 따르면 평균소비성향은 현재소득에 대한 항상소득의 비율에 의존한다.
  \item 케인즈 소비함수에서 소득이 증가할 때 평균소비성향은 항상 일정하다.
\end{choices}
\end{question}
\answer{4}
